\section*{Bab \ref{ch:g12:equilibrium} --- Kesetimbangan Kimia: Kuosien dan Tetapan}

\begin{solution}{exo:g12:equilibrium:1}
(a) $Q = \dfrac{[\ce{NH3}]^2 (c^\circ)^2}{[\ce{N2}][\ce{H2}]^3}$;
(b) $Q = \dfrac{[\ce{CH3COO-}][\ce{H3O+}]}{[\ce{CH3COOH}]\,c^\circ}$ (air,
sebagai \omterm{def:g6:solutions-and-solubility:solute}{pelarut}, tidak dituliskan);
(c) $Q = [\ce{Ca^{2+}}][\ce{CO3^{2-}}]/(c^\circ)^2$ (padatannya tidak
dituliskan);
(d) $Q = \dfrac{[\ce{FeSCN^{2+}}]\,c^\circ}{[\ce{Fe^{3+}}][\ce{SCN-}]}$.
\end{solution}

\begin{solution}{exo:g12:equilibrium:2}
$\tau = 0.020/0.050 = 0.40$: tidak tuntas.
\end{solution}

\begin{solution}{exo:g12:equilibrium:3}
$Q = 2 < K$: maju. $Q = 50 > K$: balik. $Q = K$: tidak ada perubahan.
\end{solution}

\begin{solution}{exo:g12:equilibrium:4}
Komposisinya tetap, tetapi reaksi maju dan reaksi balik terus
berlangsung dengan laju yang sama: \omterm{def:g7:atoms-and-molecules:molecule}{molekul-molekul} terus bereaksi ke
kedua arah.
\end{solution}

\begin{solution}{exo:g12:equilibrium:5}
Tidak, $K$ tidak berubah (\omterm{def:g10:reaction-progress-table:system}{keadaan akhirnya} sama); waktu untuk mencapai
kesetimbangan menjadi lebih singkat.
\end{solution}

\begin{solution}{exo:g12:equilibrium:6}
$\dfrac{x^2}{(2-x)^2} = 4$, jadi $\dfrac{x}{2-x} = 2$ dan
$x = \qty{1.33}{mol}$: \qty{0.67}{mol} asam dan \omterm{def:g11:functional-groups:alcohol}{alkohol},
\qty{1.33}{mol} \omterm{def:g11:functional-groups:carboxylic-acid}{ester} dan air (lagi-lagi dua pertiga).
\end{solution}

\begin{solution}{exo:g12:equilibrium:7}
$Q = (0.5 \times 0.5)/(0.5 \times 0.5) = 1 < 4$: maju, lebih banyak
\omterm{def:g11:functional-groups:carboxylic-acid}{ester} terbentuk.
\end{solution}

\begin{solution}{exo:g12:equilibrium:8}
Sekitar \qty{0.52}{mol} dari sisi asam dan \qty{0.79}{mol} dari sisi
\omterm{def:g11:functional-groups:carboxylic-acid}{ester}. Setelah sekitar \qty{100}{h} kedua kurva berada pada garis
putus-putus: kesetimbangan.
\end{solution}

\begin{solution}{exo:g12:equilibrium:9}
Dengan $y$ \omterm{def:g10:the-mole:amount}{mol} yang terhidrolisis: $\dfrac{y^2}{(1-y)^2} = \dfrac{1}{4}$,
jadi $\dfrac{y}{1-y} = \dfrac{1}{2}$, $y = \frac{1}{3}$: $\frac{2}{3}$
\omterm{def:g10:the-mole:amount}{mol} \omterm{def:g11:functional-groups:carboxylic-acid}{ester} dan air, $\frac{1}{3}$ \omterm{def:g10:the-mole:amount}{mol} asam dan \omterm{def:g11:functional-groups:alcohol}{alkohol}. \omterm{def:g6:pure-substances-and-mixtures:mixture}{Campuran} akhirnya
sama dengan \omterm{def:g6:pure-substances-and-mixtures:mixture}{campuran} esterifikasi.
\end{solution}

\begin{solution}{exo:g12:equilibrium:10}
Padatan adalah fase murni, yang ``konsentrasi''-nya tidak berubah:
padatan tidak dituliskan dalam $Q$. Menambahkan padatan tidak mengubah
$Q$, sehingga tidak menggeser kesetimbangan (selama masih ada padatan).
\end{solution}

\begin{solution}{exo:g12:equilibrium:11}
$K$ menetapkan perbandingan karbon dioksida dalam gas terhadap karbon
dioksida dalam air. Setelah dibuka, gas di atas air keluar ke \omterm{def:g4:air-a-mixture-of-gases:air}{udara}:
karbon dioksida dalam gas berkurang, $Q < K$, dan lebih banyak gas
meninggalkan air, sampai hampir tidak ada yang tersisa.
\end{solution}

\begin{solution}{exo:g12:equilibrium:12}
Pada kesetimbangan: asam dan \omterm{def:g11:functional-groups:alcohol}{alkohol} $\frac{1}{3}$, \omterm{def:g11:functional-groups:carboxylic-acid}{ester} dan air
$\frac{2}{3}$. Mengambil setengah dari airnya menyisakan $\frac{1}{3}$:
$Q = \dfrac{(2/3)(1/3)}{(1/3)(1/3)} = 2 < 4$, maju. Dengan $y$ tambahan
yang bereaksi: $(2/3 + y)(1/3 + y) = 4(1/3 - y)^2$, yaitu
$27y^2 - 33y + 2 = 0$, $y = 0.064$: \omterm{def:g11:functional-groups:carboxylic-acid}{esternya} naik menjadi
\qty{0.73}{mol}.
\end{solution}

\begin{solution}{exo:g12:equilibrium:13}
\omterm{def:g12:equilibrium:quotient}{Kuosien reaksi} kebalikannya memiliki pembilang dan penyebut yang
bertukar tempat: $Q' = 1/Q$, sehingga pada kesetimbangan $K' = 1/K$.
Hidrolisis: $1/4 = 0.25$.
\end{solution}

\begin{solution}{exo:g12:equilibrium:14}
$Q = (2 \times 2)/(1 \times 1) = 4 = K$: \omterm{def:g6:pure-substances-and-mixtures:mixture}{campuran} itu sudah berada dalam
kesetimbangan; komposisinya tidak berubah.
\end{solution}

\begin{solution}{exo:g12:equilibrium:15}
$n - 0.95 = 0.9025 / (4 \times 0.05) = 4.51$, jadi $n = \qty{5.5}{mol}$
\omterm{def:g11:functional-groups:alcohol}{alkohol} per \omterm{def:g10:the-mole:amount}{mol} asam.
\end{solution}

\begin{solution}{pb:g12:equilibrium:1}
\textbf{1.} \ce{CH3-COOH + CH3-CH2-OH <=> CH3-COO-CH2-CH3 + H2O}.

\textbf{2.} Reaksinya tidak tuntas: reaksi itu berhenti pada sebuah
kesetimbangan tempat kedua arah tetap berlangsung.

\textbf{3.} $Q = \dfrac{n_{\text{ester}}\, n_{\text{air}}}{n_{\text{asam}}\, n_{\text{alkohol}}}$.

\textbf{4.} $x_{\max} = \qty{1}{mol}$; $x_f = \frac{2}{3}$ \omterm{def:g10:the-mole:amount}{mol}.

\textbf{5.} $\tau = 0.67$.

\textbf{6.} Asam dan \omterm{def:g11:functional-groups:alcohol}{alkohol} masing-masing $\frac{1}{3}$ \omterm{def:g10:the-mole:amount}{mol}, \omterm{def:g11:functional-groups:carboxylic-acid}{ester} dan
air masing-masing $\frac{2}{3}$ \omterm{def:g10:the-mole:amount}{mol}.

\textbf{7.} $K = \dfrac{(2/3)^2}{(1/3)^2} = 4$.

\textbf{8.} Sepertiga terhidrolisis: \omterm{def:g11:functional-groups:carboxylic-acid}{ester} dan air $\frac{2}{3}$, asam
dan \omterm{def:g11:functional-groups:alcohol}{alkohol} $\frac{1}{3}$: \omterm{def:g6:pure-substances-and-mixtures:mixture}{campuran} yang sama, $Q_{eq} = 4$.

\textbf{9.} Reaksinya sama pada suhu yang sama, dan $K$ hanya
bergantung pada suhu.

\textbf{10.} Asam $1 - x$; \omterm{def:g11:functional-groups:alcohol}{alkohol} $3 - x$; \omterm{def:g11:functional-groups:carboxylic-acid}{ester} $x$; air $x$.

\textbf{11.} $x^2 = 4(1 - x)(3 - x) = 4(3 - 4x + x^2)$, jadi
$3x^2 - 16x + 12 = 0$.

\textbf{12.} $x = (16 \pm \sqrt{256 - 144})/6 = (16 \pm 10.58)/6$: 0.903
atau 4.43; hanya $x_f = \qty{0.903}{mol}$ yang terletak antara 0 dan 1.

\textbf{13.} \qty{90}{\percent}.

\textbf{14.} $M = \qty{88.0}{g/mol}$: $0.903 \times 88.0 = \qty{79.5}{g}$.

\textbf{15.} Setiap pengambilan air memperkecil pembilang $Q$, yang
turun di bawah $K$: sistem bergerak maju lagi.

\textbf{16.} Reaksinya akan terus berlangsung sampai asamnya habis:
\qty{100}{\percent}.

\textbf{17.} Etanol lebih murah, dan kelebihannya mudah dipisahkan dari
\omterm{def:g11:functional-groups:carboxylic-acid}{ester} (dan didaur ulang).

\textbf{18.} Hanya waktunya: reaksi yang hampir tidak melepaskan panas
memiliki tetapan yang hampir tidak berubah dengan suhu, dan pemanasan
hanya membuatnya mencapai kesetimbangan lebih cepat.

\textbf{19.} Sekitar \qty{90}{\percent} asamnya.
\end{solution}
